\section{良模空间的存在性}\label{src:7.10}
\begin{epigraph}
我认为 GIT 并未实现早期的所有期待. 最初成功构造了曲线, 阿贝尔簇, 向量丛与层的模空间, 但这些成功并没有继续推进得更远. 例如, GIT 始终未能给出紧化一般型曲面模空间的良好方法, 也没有解决高维簇的模问题.
\epigraphauthor{János Kollár \cite[第~112--113 页]{Kol06}}
\end{epigraph}
我们证明特征 $0$ 时分离良模空间存在的充要条件.
\begin{theorem}[良模空间存在性定理]\label{src:7.10.1}
设 $\mathcal X$ 是特征 $0$ 的代数闭域 $\Bbbk$ 上有限型, 具有仿射对角态射的代数叠. 存在良模空间 $\pi:\mathcal X\to X$, 其中 $X$ 是分离代数空间, 当且仅当 $\mathcal X$ 是 $\Theta$-完备的 (定义~\bookref{7.9.10}) 且 S-完备的 (定义~\bookref{7.9.14}). 此外, $X$ 固有当且仅当 $\mathcal X$ 满足固有性赋值判据的存在性部分.
\end{theorem}
\subsection{构造良模空间的策略}\label{src:7.10.1-subsection}
代数叠的局部结构定理~\bookref{7.7.1} 给出构造良模空间 $X$ 的自然策略. 对每个闭点 $x\in\mathcal X$, 有平展商展示
\[\begin{tikzcd}[ampersand replacement=\&]\mathcal W={[\Spec A/G_x]}\arrow[r,"{f}"]\arrow[d]\&\mathcal X\\W=\Spec A^{G_x}\&\end{tikzcd}\]
其中 $f$ 仿射且平展, 存在 $x$ 的原像 $w\in\mathcal W$, 使 $f$ 在 $w$ 处诱导稳定子群同构. 我们希望这些 GIT 商 $W=\Spec A^{G_x}$ 随闭点 $x$ 变化给出平展局部模型, 并可粘合成 $\mathcal X$ 的良模空间. 为此需在 $W$ 上构造平展等价关系. 因 $f$ 仿射, $\mathcal R:=\mathcal W\times_{\mathcal X}\mathcal W$ 同构于商叠 $[\Spec B/G_x]$, 且有图
\[\begin{tikzcd}[ampersand replacement=\&]\mathcal R\arrow[r,"{p_1}"]\arrow[d]\&\mathcal W\arrow[r,"{f}"]\arrow[d]\&\mathcal X\\R\arrow[r,"{q_1}"]\&W\&\end{tikzcd}\]
其中 $R=\Spec B^{G_x}$. 若 $q_1,q_2:R\rightrightarrows W$ 定义平展等价关系, 则代数空间商 $W/R$ 是 $f(\mathcal W)\subseteq\mathcal X$ 的良模空间的候选. Luna 基本引理~\bookref{7.5.33} 给出 $q_1,q_2$ 平展的条件: 对所有闭点 $r\in\mathcal R$, 要求 (a) $p_1(r),p_2(r)\in\mathcal W$ 都是闭点; (b) $p_1,p_2$ 在 $r$ 处诱导稳定子群同构.

另一方面, 已知 $f(w)$ 闭, 且 $f$ 在给定原像 $w$ 处诱导稳定子群同构. 我们希望找到 $w$ 的开邻域 $\mathcal U$, 使 $f|_{\mathcal U}$ 满足 (a) 将闭点映为闭点; (b) 在闭点处诱导稳定子群同构, 且这些条件对基变换稳定.\footnote{由 Luna 基本引理~\bookref{7.5.30}, 确实存在 $\rho=(w,w,\mathrm{id})\in\mathcal R$ 的开邻域 $\mathcal R_1$ (分别地, $\mathcal R_2$), 使 $p_1$ (分别地, $p_2$) 满足 (a),(b); 但不清楚这些轨迹是否 $\mathcal R$-不变, 即是否为 $\mathcal W$ 的开子叠的原像.} 条件 (b) 对基变换稳定, 条件 (a) 则不然. 下文引入更强的 $\Theta$-满射性 (定义~\bookref{7.10.8}), 它对基变换稳定且蕴含 (a).

在构造中, $\mathcal X$ 的 $\Theta$-完备性保证 $w$ 的某个开邻域上的 $\Theta$-满射性, 从而保证 (a) 及其所有基变换 (命题~\bookref{7.10.13}); S-完备性则保证 $w$ 的某个开邻域上的 (b) (命题~\bookref{7.10.20}).
\subsection{反例}\label{src:7.10.2-subsection}
以下例子没有良模空间. 上述证明策略在这些例子中失效, 有助于理解两种完备性的作用. 在代数闭域 $\Bbbk$ 上工作. 前三个例子 $[\mathbb P^1/\mathbb G_m]$, $[C/\mathbb G_m]$ 与 $[(\mathbb A^2\setminus0)/\mathbb G_m]$ 不满足 $\Theta$-完备性, 见例~\bookref{7.9.28}.
\begin{example}\label{src:7.10.2}
商叠 $[\mathbb P^1/\mathbb G_m]$ 没有良模空间, 其中作用为 $t\cdot[x:y]=[tx:y]$. 对具有良模空间的叠, 每个 $\Bbbk$-点的闭包中有唯一闭点 (定理~\bookref{7.5.11}(2)). 此处失败, 因为 $[1:1]$ 特化为两个闭点 $[1:0],[0:1]$. 另一论证是: 若存在良模空间, 它必为 $[\mathbb P^1/\mathbb G_m]\to\Spec\Bbbk$ (此映射对到代数空间的映射具有泛性质), 但 Serre 仿射性判据~\bookref{5.5.16} 将蕴含复合 $\mathbb P^1\to[\mathbb P^1/\mathbb G_m]\to\Spec\Bbbk$ 仿射, 矛盾.

有两个同构于 $\mathbb A^1/\mathbb G_m$ 的开子叠 $\mathcal U_1,\mathcal U_2$, 各有良模空间 $\pi_i:\mathcal U_i\to\Spec\Bbbk$, 却不能粘合: 交 $\mathcal U_1\cap\mathcal U_2$ 是两者中的开点, 并非 $\pi_i$ 下开子概形的原像. 平展展示 $f:\mathcal W:=\mathcal U_1\coprod\mathcal U_2\to[\mathbb P^1/\mathbb G_m]$ 满足 (a),(b), 但其基变换 $p_1,p_2:\mathcal W\times_{[\mathbb P^1/\mathbb G_m]}\mathcal W=\mathcal U_1\coprod\mathcal U_2\coprod\mathcal U_1\cap\mathcal U_2\to\mathcal W$ 不满足 (b): 交中的闭点在任一投影下均映到非闭点.
\end{example}
\begin{example}\label{src:7.10.3}
设 $C$ 是带 $\mathbb G_m$-作用的射影结点三次曲线. 商 $[C/\mathbb G_m]$ 有两个点, 一个开点, 一个闭点; 虽无上述拓扑障碍, 仍因 $C$ 射影而非仿射而没有良模空间. 将 $C$ 视为两个 $\mathbb P^1$ 在 $0,\infty$ 处相接得到的结点并 $X'$ 对 $\mathbb Z/2$ 旋转作用的商, 得到有限平展覆盖 $[X'/\mathbb G_m]\to[X/\mathbb G_m]$. 删去一个原点, 得到仿射平展覆盖 $\mathcal W=[\Spec(\Bbbk[x,y]/xy)/\mathbb G_m]\to[C/\mathbb G_m]$, 其中 $t\cdot(x,y)=(tx,t^{-1}y)$. 此映射将闭点映为闭点, 但投影 $\mathcal W\times_{[C/\mathbb G_m]}\mathcal W\rightrightarrows\mathcal W$ 不然.
\end{example}
\begin{example}\label{src:7.10.4}
令 $\mathbb G_m$ 通过 $t\cdot(x,y)=(tx,y)$ 作用于 $\mathbb A^2$, 置 $\mathcal X=[\mathbb A^2/\mathbb G_m]\setminus0$. 点 $p=(1,0)$ 闭且稳定子平凡. 将 $z$ 送到 $(z,1)$ 的开浸入 $f:\mathbb A^1\hookrightarrow\mathcal X$ 是平展商展示. $f(0)$ 闭, 但 $z\ne0$ 时 $f(z)$ 不闭. 由 $(x,y)\mapsto y$ 定义的 $\mathcal X\to\mathbb A^1$ 非良模空间, 因 $\mathbb A^2\setminus0$ 非仿射.
\end{example}
不满足 $\Theta$-完备性的例子自然出现在模问题中, 例如从 $\cBun_{r,d}(C)^{\mathrm{ss}}$ 中删去一个多重稳定但非稳定的向量丛. 以下三个例子不满足 S-完备性.
\begin{example}\label{src:7.10.5}
设 $\operatorname{char}(\Bbbk)\ne2$, 令 $G=\mathbb Z/2$ 通过交换两份仿射直线作用于非分离并 $U=\mathbb A^1\cup_{x\ne0}\mathbb A^1$. 商叠 $[U/G]$ 除原点外处处有 $\mathbb Z/2$ 稳定子. 它是惯性拟有限但不有限的 Deligne--Mumford 叠, 非 S-完备. 习题~\bookref{5.4.16} 已用它说明 Keel--Mori 定理~\bookref{5.4.6} 中有限惯性假设的必要性. 预复合一份 $\mathbb A^1$ 的包含映射, 得到仿射平展态射 $\mathbb A^1\to[U/G]$, 它在 $0$ 处保持稳定子, 却不在 $0$ 的任何开邻域上保持稳定子.
\end{example}
\begin{example}\label{src:7.10.6}
$\mathbb P^1$ 上四个无序点的模叠 $\mathcal X=[\operatorname{Sym}^4\mathbb P^1/\mathrm{PGL}_2]$ 的 Deligne--Mumford 轨迹 $\mathcal X^{\mathrm{DM}}$ 非分离 (例~\bookref{5.4.17}), 故非 S-完备. 但由四个互异点组成的稳定轨迹 $\mathcal X^{\mathrm{s}}$ 分离, 半稳定轨迹 $\mathcal X^{\mathrm{ss}}=\mathcal X^{\mathrm{DM}}\cup\{[0:0:\infty:\infty]\}$ 有射影良模空间.
\end{example}
\begin{example}\label{src:7.10.7}
考虑 $G=\mathbb G_m\times\mathbb Z/2$ 在 $X=\mathbb A^2\setminus0$ 上的作用 $t\cdot(a,b)=(ta,t^{-1}b)$ 与 $-1\cdot(a,b)=(b,a)$. 每个满足 $ab\ne0$ 的点 $(a,b)$ 都被二阶元素 $(a/b,-1)\in G$ 固定. 商叠 $[X/G]$ 是非分离 Deligne--Mumford 叠, 非 S-完备且没有良模空间; 但 $[\mathbb A^2/G]\to\Spec\Bbbk[xy]$ 是良模空间.
\end{example}
\subsection{\texorpdfstring{$\Theta$}{Θ}-完备性与 \texorpdfstring{$\Theta$}{Θ}-满射性}\label{src:7.10.3-subsection}
将闭点映为闭点的性质对基变换不稳定 (例~\bookref{7.10.2}, \bookref{7.10.3}). 引入更强且性质更好的 $\Theta$-满射性. 主要结果是: 只要 $\mathcal X$ 是 $\Theta$-完备的, 平展商展示 $([\Spec A/G_x],w)\to(\mathcal X,x)$ 就在 $w$ 的某个开邻域上 $\Theta$-满射 (命题~\bookref{7.10.13}). 如第~\bookref{7.10.1-subsection} 节所述, 这对存在性定理~\bookref{7.10.1} 的证明至关重要.
\begin{definition}\label{src:7.10.8}
设 $f:\mathcal X\to\mathcal Y$ 是代数叠的态射, $x\in\mathcal X(\Bbbk)$ 是几何点. 若每个图
\begin{align}\label{src:7.10.9}
\begin{tikzcd}[ampersand replacement=\&]\Spec\Bbbk\arrow[r,"{x}"]\arrow[d,"{1}"']\&\mathcal X\arrow[d,"{f}"]\\\Theta_{\Bbbk}\arrow[r]\arrow[ur,dashed]\&\mathcal Y\end{tikzcd}
\end{align}
都有提升, 则称 $f$ 在 $x$ 处 \newterm{$\Theta$-满射}{$\Theta$-surjective}. 若在每个几何点处如此, 则称 $f$ 为 $\Theta$-满射.
\end{definition}
此概念显然对基变换稳定. 若Noether代数叠 $\mathcal Y$ 的对角态射仿射且拟有限, 则每个态射 $f:\mathcal X\to\mathcal Y$ 都是 $\Theta$-满射的, 因为每个 $\Theta_{\Bbbk}\to\mathcal Y$ 都经由 $\Spec\Bbbk$ 分解 (引理~\bookref{7.9.21}). 下述引理给出提升唯一及定义与几何点选择无关的条件.
\begin{lemma}\label{src:7.10.10}
设 $f:\mathcal X\to\mathcal Y$ 是Noether代数叠之间分离, 可表且有限型的态射.
\begin{enumerate}[(1)]
\item 图~\bookeqref{7.10.9} 的每个提升唯一.
\item 若 $f$ 在 $x\in\mathcal X(\Bbbk)$ 处 $\Theta$-满射, 则在表示 $|\mathcal X|$ 中同一点的每个其他几何点 $x'\in\mathcal X(\Bbbk')$ 处也 $\Theta$-满射.
\end{enumerate}
\end{lemma}
\begin{proof}
(1) 来自下降与分离性赋值判据. 为证明 (2), 只需证明对代数闭域扩张 $\Bbbk\to\Bbbk'$, 提升 $\Theta_{\Bbbk'}\to\mathcal X$ 的存在蕴含 $\Theta_{\Bbbk}\to\mathcal X$ 的存在. 写 $\Bbbk'=\bigcup_\lambda A_\lambda$ 为有限生成 $\Bbbk$-子代数的并. 由极限方法 (第~\bookref{B.3} 节), 存在 $\Spec A_\lambda\to\mathcal X$ 的提升 $\Theta_{A_\lambda}\to\mathcal X$. 沿 $\Spec A_\lambda$ 的闭点限制, 得到 $\Bbbk$ 上的提升.
\end{proof}
\begin{proposition}\label{src:7.10.11}
设 $f:\mathcal X\to\mathcal Y$ 是代数闭域 $\Bbbk$ 上有限型, 具有仿射对角态射的代数叠之间的态射. 假设 $\mathcal Y$ 的闭点有线性约化稳定子. 若 $f$ 为 $\Theta$-满射, 则 $f$ 将闭点映为闭点.
\end{proposition}
\begin{proof}
设 $x\in|\mathcal X|$ 闭, $f(x)\rightsquigarrow y_0$ 是到闭点的特化. 由失稳定理~\bookref{8.3.8}, 此特化可由 $\Theta\to\mathcal Y$ 实现. $\Theta$-满射性将它提升为 $g:\Theta\to\mathcal X$, 满足 $g(1)=x$. 因 $x$ 闭, 此提升只能对应于平凡特化 $x\rightsquigarrow x$. 所以 $f(x)=y_0$ 闭.
\end{proof}
\begin{remark}\label{src:7.10.12}
逆命题不成立. 例~\bookref{7.10.3} 中, $C$ 是带 $\mathbb G_m$-作用的结点三次曲线. 平展态射
\[[\Spec(\Bbbk[x,y]/(xy))/\mathbb G_m]\to[C/\mathbb G_m]\]
将闭点映为闭点, 却非 $\Theta$-满射.
\end{remark}
\begin{proposition}\label{src:7.10.13}
设 $\mathcal X$ 是代数闭域 $\Bbbk$ 上有限型, 具有仿射对角态射的代数叠, 其闭点有线性约化稳定子. 设 $x\in|\mathcal X|$ 闭, $f:([\Spec A/G_x],w)\to(\mathcal X,x)$ 是在 $w$ 处诱导稳定子群同构的仿射平展态射, 令 $\pi:[\Spec A/G_x]\to\Spec A^{G_x}$. 若 $\mathcal X$ 是 $\Theta$-完备的, 则存在 $\pi(w)$ 的仿射开邻域 $U\subseteq\Spec A^{G_x}$, 使 $f|_{\pi^{-1}(U)}:\pi^{-1}(U)\to\mathcal X$ 为 $\Theta$-满射.
\end{proposition}
\begin{proof}
令 $\mathcal W=[\Spec A/G_x]$, 定义 $\Sigma_f\subseteq|\mathcal W|$ 为使 $f$ 在该点处 $\Theta$-满射的点集. 先在 $\mathcal X\cong[\Spec B/G]$, $G$ 线性约化时证明 $\Sigma_f$ 开. Zariski 主定理~\bookref{7.1.17} 给出分解 $f:\mathcal W\xrightarrow{j}\widetilde{\mathcal X}\xrightarrow{\nu}\mathcal X$, 其中 $j$ 为开浸入, $\nu$ 有限. 由引理~\bookref{7.9.19}, $[\Spec B/G]$ 是 $\Theta$-完备的. 由命题~\bookref{7.9.20}, 因 $\nu$ 仿射, $\widetilde{\mathcal X}$ 亦然. $\nu$ 有限蕴含 $\Sigma_j=\Sigma_f$, 故可假设 $f$ 是开浸入. 令 $\mathcal Z$ 是 $\mathcal W$ 在 $\mathcal X$ 中的约化补, $\pi:\mathcal X\to\Spec B^G$ 为良模空间. 断言 $|\mathcal W|\setminus\Sigma_f=\pi^{-1}(\pi(|\mathcal Z|))\cap|\mathcal W|$. 包含关系 ``$\subseteq$'' 显然: $\mathcal X\setminus\pi^{-1}(\pi(|\mathcal Z|))\hookrightarrow\mathcal X$ 是代数空间的 $\Theta$-满射态射 $X\setminus\pi(|\mathcal Z|)\hookrightarrow X$ 的基变换. 反之, 设 $y$ 在右侧, 由几何点 $\Spec L\to\mathcal X$ 表示. 令 $z\in|\mathcal Z_L|$ 为 $y\in|\Theta_L|$ 的闭包中的唯一闭点, 取表示特化 $y\rightsquigarrow z$ 的态射 $\Theta_L\to\mathcal X_L$ (推论~\bookref{8.3.8}). 此态射不能提升到 $\mathcal W$, 故 $y\notin\Sigma_f$.

现在断言 $\Sigma_f$ 可构造. 用局部结构定理~\bookref{7.7.1} 选择仿射平展满射 $g:\mathcal X'=[\Spec B/G]\to\mathcal X$, 其中 $G$ 线性约化. 令 $\mathcal W'=\mathcal W\times_{\mathcal X}\mathcal X'$, 投影为 $g':\mathcal W'\to\mathcal W$, $f':\mathcal W'\to\mathcal X'$. 已知 $\Sigma_{f'}$ 开, 由 Chevalley 定理~\bookref{4.3.31}, 只需证明 $\mathcal W\setminus\Sigma_f=g'(\mathcal W'\setminus\Sigma_{f'})$. 为此只需证明每个 $h:\Theta_L\to\mathcal X$ ($L$ 代数闭) 都提升为 $h':\Theta_L\to\mathcal X'$. 取 $h(0)$ 的原像 $x'\in\mathcal X'(L)$. 因 $g$ 可表且平展, 稳定子映射 $G_{x'}\to G_{h(0)}$ 单射且余核有限. 所以 $h$ 诱导的 $\mathbb G_{m,L}\to G_{h(0)}$ 经由 $G_{x'}$ 分解, 可将 $h|_{B\mathbb G_{m,L}}$ 提升到 $\mathcal X'$. 对 $B\mathbb G_{m,L}\hookrightarrow\Theta_L$ 的第 $n$ 个幂零加厚 $\mathcal X_n$, 变形理论 (命题~\bookref{7.7.12}) 给出相容提升 $\mathcal X_n\to\mathcal X'$, 凝聚 Tannaka 对偶性~\bookref{7.6.9} 将其延拓为 $\Theta_L\to\mathcal X'$.

因 $\Sigma_f$ 可构造且 $w\in\Sigma_f$, 为证明它开, 只需证明 $w$ 的每个泛化 $\xi\rightsquigarrow w$ 均在 $\Sigma_f$ 中. 取表示此特化的完备 DVR 态射 $h:\Spec R\to\mathcal W$, 分式域, 剩余域分别为 $K,\kappa$. 断言存在提升 (因 $f$ 分离而必唯一)
\begin{align}\label{src:7.10.14}
\begin{tikzcd}[ampersand replacement=\&]\Spec K\arrow[r,"{h_K}"]\arrow[d,"{1}"']\&\mathcal W\arrow[d,"{f}"]\\\Theta_K\arrow[r,"{g}"]\arrow[ur,dashed,"{\widetilde g}"']\&\mathcal X\end{tikzcd}
\end{align}
这便说明 $\xi\in\Sigma_f$. 先由 $\Theta$-完备性构造提升
\[\begin{tikzcd}[ampersand replacement=\&]\Theta_R\setminus0\arrow[r]\arrow[d]\&\mathcal X\\\Theta_R\arrow[ur,dashed,"{q}"']\&\end{tikzcd}\]
因 $\mathcal W\to\mathcal X$ 在 $w$ 处保持稳定子, $q|_{B\mathbb G_{m,\kappa}}$ 有提升 $B\mathbb G_{m,\kappa}\to\mathcal W$. $\Theta_R$ 沿 $B\mathbb G_{m,\kappa}$ 凝聚完备 (定理~\bookref{7.6.13}), 因而由变形理论~\bookref{7.7.12} 与凝聚 Tannaka 对偶性~\bookref{7.6.9} 可构造提升
\[\begin{tikzcd}[ampersand replacement=\&]B\mathbb G_{m,\kappa}\arrow[r]\arrow[d,hook]\&\mathcal W\arrow[d,"{f}"]\\\Theta_R\arrow[r,"{q}"]\arrow[ur,dashed,"{\widetilde q}"']\&\mathcal X\end{tikzcd}\]
$\widetilde q|_{\Spec R}$ 与 $h$ 在闭点处相同, 而 $f$ 平展, 故两者 $2$-同构. 所以 $\widetilde g:=\widetilde q|_{\Theta_K}$ 是图~\bookeqref{7.10.14} 的提升.
\end{proof}
$\Theta$-完备叠的 $\Bbbk$-点的拓扑类似于 GIT 商叠的拓扑.
\begin{exercise}\label{src:7.10.15}
设 $\mathcal X$ 是代数闭域 $\Bbbk$ 上有限型, 具有仿射对角态射的代数叠. 假设它 $\Theta$-完备, 且闭点有线性约化稳定子. 证明每个 $\Bbbk$-点的闭包含唯一闭点.
\end{exercise}
\begin{exercise}\label{src:7.10.16}
在习题~\bookref{7.10.15} 的假设下, 若 $\mathcal X$ 还有唯一闭点, 证明 $\mathcal X\cong[\Spec A/G_x]$, 且 $A^{G_x}$ 是剩余域为 $\Bbbk$ 的阿廷局部 $\Bbbk$-代数.
\end{exercise}
\subsection{无穿孔惯性}\label{src:7.10.4-subsection}
引入代数叠的无穿孔惯性, 并证明它蕴含每个平展商展示 $f:([\Spec A/G_x],w)\to(\mathcal X,x)$ 在 $w$ 的某个开邻域上保持稳定子 (命题~\bookref{7.10.20}).
\begin{definition}\label{src:7.10.17}
若对每个闭点 $x\in|X|$ 及每个形式通用态射 $p:(T,t)\to(\mathcal X,x)$, 其中 $T$ 是闭点为 $t$ 的局部环的谱, 惯性群概形 $\underline{\operatorname{Aut}}_{\mathcal X}(p)\to T$ 的每个连通分支都与 $t$ 上的纤维非空相交, 则称Noether代数叠具有\newterm{无穿孔惯性}{unpunctured inertia}.
\end{definition}
\begin{remark}\label{src:7.10.18}
此处 $(T,t)\to(\mathcal X,x)$ 形式通用, 指完备化的映射 $\widehat T\to\mathcal X$ 在 $t$ 处按定义~\bookref{C.4.2} 形式通用.
\end{remark}
\begin{remark}\label{src:7.10.19}
无穿孔性与 \cite[第~3.3 节]{GR71} 定义的 $\underline{\operatorname{Aut}}_{\mathcal X}(p)\to T$ 的纯性有关 (另见 \cite[Tag 0CV5]{SP}). 若 $T$ 为严格 Hensel 局部环的谱, 纯性要求: 对任意 $s\in T$ 及纤维 $\underline{\operatorname{Aut}}_{\mathcal X}(p)_s$ 的伴随点 $\gamma$, $\gamma$ 在 $\underline{\operatorname{Aut}}_{\mathcal X}(p)$ 中的闭包与 $T$ 的闭点 $t$ 上的纤维非空相交.
\end{remark}
\begin{proposition}\label{src:7.10.20}
设 $\mathcal X$ 是代数闭域 $\Bbbk$ 上有限型, 具有仿射对角态射的代数叠, $x$ 是具有线性约化稳定子 $G_x$ 的闭点. 设 $f:([\Spec A/G_x],w)\to(\mathcal X,x)$ 仿射且平展, 在 $w$ 处诱导稳定子群同构, 令 $\pi:[\Spec A/G_x]\to\Spec A^{G_x}$. 若 $\mathcal X$ 具有无穿孔惯性, 则存在 $\pi(w)$ 的仿射开邻域 $U$, 使 $f|_{\pi^{-1}(U)}$ 在所有点处诱导稳定子群同构.
\end{proposition}
\begin{proof}
置 $\mathcal W=[\Spec A/G_x]$. 只需找到 $w$ 的开邻域 $\mathcal U$, 使 $f|_{\mathcal U}$ 诱导同构 $I_{\mathcal U}\to\mathcal U\times_{\mathcal X}I_{\mathcal X}$. 考虑笛卡尔图 (习题~\bookref{4.2.15})
\[\begin{tikzcd}[ampersand replacement=\&]I_{\mathcal W}\arrow[r]\arrow[d]\&\mathcal W\times_{\mathcal X}I_{\mathcal X}\arrow[d]\\\mathcal W\arrow[r]\&\mathcal W\times_{\mathcal X}\mathcal W\end{tikzcd}\]
因 $f$ 分离且平展, 上方态射有限且平展. 令 $\mathcal Z\subseteq\mathcal W\times_{\mathcal X}I_{\mathcal X}$ 为其不是同构的开闭子叠. 因 $f$ 在 $w$ 处保持稳定子, $w\notin p_1(\mathcal Z)$.

考虑来自闭点为 $t$ 的局部环的谱的形式光滑态射 $(T,t)\to(\mathcal X,x)$. 无穿孔惯性蕴含 $\mathcal Z$ 在 $\mathcal W\times_{\mathcal X}I_{\mathcal X}\times_{\mathcal X}T$ 中的原像为空: 否则原像的一个非空连通分支必与 $t$ 上的纤维非空相交, 与 $w\notin p_1(\mathcal Z)$ 矛盾. 因而 $w\notin\overline{p_1(\mathcal Z)}$. 取 $\mathcal U=\mathcal W\setminus\overline{p_1(\mathcal Z)}$, 所诱导的惯性态射便是同构.
\end{proof}
\begin{proposition}\label{src:7.10.21}
设 $\mathcal X$ 是Noether代数叠.
\begin{enumerate}[(1)]
\item 若 $\mathcal X$ 具有拟有限惯性, 则它具有无穿孔惯性当且仅当惯性有限.
\item 若 $\mathcal X$ 的稳定子群连通, 则它具有无穿孔惯性.
\end{enumerate}
\end{proposition}
\begin{proof}
惯性有限时, $\underline{\operatorname{Aut}}_{\mathcal X}(p)\to T$ 有限, 每个连通分支的像显然含唯一闭点 $t$. 反之可假设 $T$ 为 Hensel 局部环的谱, 此时 $\underline{\operatorname{Aut}}_{\mathcal X}(p)=G\coprod H$, $G\to T$ 有限, $H$ 在 $t$ 上的纤维为空 (命题~\bookref{B.5.9}). 若 $T$ 非空 (即该自同构群概形不有限), 则 $T$ 的任一连通分支不与中心纤维相交, 所以 $\mathcal X$ 的惯性非无穿孔.

对 (2), $\underline{\operatorname{Aut}}_{\mathcal X}(p)\to U$ 的所有纤维连通, 所以每个连通分支都与包含单位截面的分支相交.
\end{proof}
\begin{remark}\label{src:7.10.22}
对于稳定子群连通的代数叠, 例如曲线上半稳定向量丛的模叠 $\cBun_{r,d}^{\mathrm{ss}}(C)$, 命题~\bookref{7.10.21}(2) 已保证无穿孔惯性. 因而证明此模叠的良模空间存在时, 不需要下述更深的定理~\bookref{7.10.23}.
\end{remark}
\subsection{S-完备性蕴含无穿孔惯性}\label{src:7.10.5-subsection}
\begin{theorem}\label{src:7.10.23}
设 $\mathcal X$ 是代数闭域 $\Bbbk$ 上有限型, 具有仿射对角态射的代数叠, 闭点的稳定子线性约化. 若 $\mathcal X$ 是 S-完备的, 则它具有无穿孔惯性.
\end{theorem}
\begin{proof}
设 $x$ 闭, $p:(U,u)\to(\mathcal X,x)$ 是来自局部环的谱的形式光滑态射, $H\subseteq\underline{\operatorname{Aut}}_{\mathcal X}(p)$ 是连通分支. 投影 $H\to U$ 的像可构造, 闭包含 $u$. 选择剩余域为 $\Bbbk$ 的 DVR $R$ 及 $\Spec R\to U$, 特殊点映到 $u$, 泛点落在 $H\to U$ 的像中. 令 $\xi:\Spec R\to U\xrightarrow{p}\mathcal X$ 为复合. 作剩余域不变的 DVR 扩张后, 可假设 $\Spec K\to U$ 提升到 $H$, 得到交换图
\[\begin{tikzcd}[ampersand replacement=\&]\Spec K\arrow[r]\arrow[d]\&\Spec R\arrow[d]\arrow[dr,"{\xi}"]\&\\H\arrow[r]\&U\arrow[r,"{p}"]\&\mathcal X\end{tikzcd}\]
令 $H_K$ 为沿 $\Spec K\to U$ 的基变换. 断言可选择有限阶的有限型点 $g\in H_K$. 从任意有限型点出发, 有限扩张 $K$ 后可作 Jordan 分解 $g=g_sg_u$, 其中 $g_s$ 半单, $g_u$ 幺幂 (第~\bookref{B.1.3-subsection} 节). 由 $g_s$ 生成的约化闭 $K$-子群 $H'\subseteq\underline{\operatorname{Aut}}_{\mathcal X}(p)_K$ 可对角化, 可将 $g_s$ 替换为其中仍与 $g_u$ 交换的有限阶元素. 若 $\operatorname{char}(K)>0$, $g_u$ 有限阶, 即得结论. 若特征为 $0$, $g_u$ 在 $G$ 的单位分支中, 所以 $g$ 与有限阶元素 $g_s$ 在同一分支上.

断言作剩余域不变的扩张后, 存在 $\xi':\Spec R\to\mathcal X$, 满足 $\xi'_K\simeq\xi_K$, 且 $g$ 延拓为 $\xi'$ 的自同构. 这将完成证明: $g$ 的闭包与 $\underline{\operatorname{Aut}}_{\mathcal X}(p)\to U$ 在 $u$ 上的纤维相交, 故 $H$ 与中心纤维相交. 若 $\mathcal X\cong[\Spec A/\mathrm{GL}_n]$, 断言正是下述命题~\bookref{7.10.24}. 用局部结构定理~\bookref{7.7.1} 化归到此情形: 取平展商展示 $f:([\Spec A/G_x],w)\to(\mathcal X,x)$, 作剩余域不变的扩张后可将 $\xi$ 提升为 $\widetilde\xi:\Spec R\to[\Spec A/G_x]$, 满足 $\widetilde\xi(0)=w$. 为证明 $g$ 提升为 $\widetilde g\in\operatorname{Aut}(\widetilde\xi_K)$, 使用 S-完备性. 沿 $g$ 将 $\xi$ 与自身粘合, 得到
\[\Spec R\cup_{\Spec K}\Spec R=\phi_R\setminus0\to\mathcal X.\]
此映射延拓为 $h:\phi_R\to\mathcal X$. 因 $\xi(0)=x$ 且 $x$ 闭, $h(0)=x$. $f$ 在 $w$ 处保持稳定子, 所以 $h|_{B\mathbb G_m}$ 提升为像为 $w$ 的 $\widetilde h_0:B\mathbb G_m\to[\Spec A/G_x]$. 变形理论~\bookref{7.7.12} 给出沿 $0$ 的幂零加厚上的相容提升, 凝聚 Tannaka 对偶性~\bookref{7.6.9} 给出图中的提升
\[\begin{tikzcd}[ampersand replacement=\&]B\mathbb G_m\arrow[r,"{\widetilde h_0}"]\arrow[d,hook]\&{[\Spec A/G_x]}\arrow[d,"{f}"]\\\phi_R\arrow[r,"{h}"]\arrow[ur,dashed,"{\widetilde h}"']\&\mathcal X\end{tikzcd}\]
因 $f$ 仿射且平展, $\widetilde h|_{s\ne0}$ 与 $\widetilde h|_{t\ne0}$ 都同构于 $\widetilde\xi$, 从而 $\widetilde h|_{\phi_R\setminus0}$ 给出 $g$ 的提升 $\widetilde g$. 最后应用命题~\bookref{7.10.24}, 构造 $\widetilde\xi':\Spec R\to[\Spec A/G_x]$, 满足 $\widetilde\xi'(0)=w$, $\widetilde\xi_K\simeq\widetilde\xi'_K$, 且 $\widetilde g$ 延拓为其自同构. 复合 $\xi'':=f\circ\widetilde\xi'':\Spec R\to\mathcal X$ 满足断言. 另见 \cite[定理~5.2]{AHLH23}.
\end{proof}
证明使用了下述商叠的赋值判据.
\begin{proposition}\label{src:7.10.24}
设 $\mathcal X=[\Spec A/G]$, 其中 $\Spec A$ 是代数闭域 $\Bbbk$ 上带线性约化群 $G$ 作用的有限型仿射概形. 设 $x\in|\mathcal X|$ 闭. 则 $\mathcal X$ 满足以下性质:

($\star$) 对每个剩余域为 $\Bbbk$, 分式域为 $K$ 的 DVR $R$, 每个满足 $\xi(0)\simeq x$ 的 $\xi:\Spec R\to\mathcal X$, 以及每个有限阶 $K$-点 $g\in\operatorname{Aut}_{\mathcal X}(\xi_K)$, 存在 DVR 扩张 $R\to R'$ (令 $K'=\operatorname{Frac}(R')$) 及 $\xi':\Spec R'\to\mathcal X$, 使 $\xi'(0)\simeq x$, $\xi'_{K'}\simeq\xi_{K'}$, 且 $g|_{K'}$ 延拓为 $\xi'$ 的自同构.
\end{proposition}
\begin{remark}\label{src:7.10.25}
换言之, 对每个 $\xi:\Spec R\to\Spec A$ 和有限阶元素 $g\in G_{\xi_K}\subseteq G(K)$, 在 DVR 扩张 $R\subseteq R'$ 后存在 $h\in G(K')$, 使 $h\cdot\xi_{K'}$ 延拓为 $\xi':\Spec R'\to\Spec A$, 满足 $\xi'(0)\in G\cdot\xi(0)$, 且 $h^{-1}g|_{K'}h$ 延拓为 $G$ 的 $R'$-点. 例如 $G=\mathbb G_m\rtimes\mathbb Z/2$ 在 $\mathbb A^2$ 上通过 $t\cdot(a,b)=(ta,t^{-1}b)$, $-1\cdot(a,b)=(b,a)$ 作用. 当 $ab\ne0$, $(a,b)$ 被二阶元素 $(a/b,-1)$ 固定. 考虑 $\xi:\Spec R=\Bbbk[z]\to\mathbb A^2$, $z\mapsto(z^2,z)$. 元素 $g=(z^{-1},-1)\in G(\Bbbk((z)))$ 稳定 $\xi_K$, 却不延拓到 $G(\Bbbk[z])$. 但可取二次分歧扩张 $\Bbbk[z]\to\Bbbk[\sqrt z]$, 定义 $\xi':\Spec\Bbbk[\sqrt z]\to\mathbb A^2$, $\sqrt z\mapsto((\sqrt z)^3,(\sqrt z)^3)$. 泛点上的同构由 $h=(\sqrt z,-1)$ 给出, 且
\[g|_{\Bbbk((\sqrt z))}=(\sqrt z,-1)^{-1}\cdot g|_{K'}\cdot(\sqrt z,-1)=(1,-1)\]
延拓为 $G$ 的 $\Bbbk[\sqrt z]$-点.
\end{remark}
\begin{proof}
选择嵌入 $G\hookrightarrow\mathrm{GL}_n$, 将商叠替换为 $[(\Spec A\times^{G_x}\mathrm{GL}_n)/\mathrm{GL}_n]$, 可假设 $G=\mathrm{GL}_n$. 先对平凡作用的商叠 $[\Spec A/G]=\Spec A\times BG$ 验证 ($\star$). 因 $R$ 局部, $G=\mathrm{GL}_n$, 复合 $\Spec R\to[\Spec A/G]\to BG$ 对应于平凡 $G$-丛. 只需证明扩张 DVR 后每个有限阶 $g\in G(K)$ 都共轭于 $G(R)$ 中的元素. 扩张后可将 $g$ 化为 Jordan 标准形. 因 $g$ 有限阶, 对角元为某个 $r$ 的 $r$ 次单位根. $r$ 次单位根群 $\boldsymbol\mu_r$ 在 $\Spec R$ 上有限, 所以标准形的各项均在 $R$ 中.

若 $\mathcal X=[X/G]$, $X$ 在 $\Bbbk$ 上固有, 则 ($\star$) 成立, 但 $\xi'(0)$ 未必同构于 $x$. 因 $p:\mathcal X\to BG$ 固有且可表, 对 DVR 态射 $\xi$ 有群概形闭浸入 $\underline{\operatorname{Aut}}_{\mathcal X}(\xi)\hookrightarrow\underline{\operatorname{Aut}}_{BG}(p\circ\xi)$. 此外, $\Spec R\to BG$ 的泛点到 $[X/G]$ 的任一提升都延拓为唯一的 $\Spec R\to BG$. 给定 $g\in\operatorname{Aut}_{\mathcal X}(\xi_K)$, 利用 $BG$ 的 ($\star$), 扩张后取 $\eta:\Spec R\to BG$, 满足 $\eta_K\simeq(p\circ\xi)_K$, 且 $g|_K$ 延拓为 $\operatorname{Aut}_{BG}(\eta)$ 的 $R$-点. 将 $\eta$ 提升为 $\xi':\Spec R\to[X/G]$, 满足 $\xi'_K\simeq\xi_K$, 则 $g|_K$ 延拓为 $\xi'$ 的自同构.

验证 $[\Spec A/G]$ 的 ($\star$) 时可假设 $A$ 约化. 将此叠视为 $\Spec A^G\times BG$ 上仿射有限型代数叠, 可选择该底上的向量丛 $\mathcal E$ 及 $A^G$ 上的 $G$-等变嵌入 $\Spec A\hookrightarrow\mathbb A(\mathcal E|_{\Spec A^G})$. 将后者视为 $\mathbb P(\mathcal E|_{\Spec A^G}\oplus\mathcal O_{\Spec A^G})$ 的开子概形, 令 $X$ 为 $\Spec A$ 在其中的闭包. 得到 $G$-等变图
\begin{align}\label{src:7.10.26}
\begin{tikzcd}[ampersand replacement=\&]\Spec A\arrow[r,hook]\arrow[dr]\&X\arrow[d]\\\&\Spec A^G\end{tikzcd}
\end{align}
其中 $X$ 是约化射影概形, 补 $X\setminus\Spec A$ 是丰沛 $G$-不变 Cartier 因子 $E$ 的支集. 断言 $\Spec A$ 恰为 $X$ 关于 $\mathcal O_X(E)$ 的半稳定轨迹 (习题~\bookref{8.2.11}). 典范不变截面 $s:\mathcal O_X\to\mathcal O_X(E)$ 在 $\Spec A$ 上为同构, 故 $\Spec A\subseteq X^{\mathrm{ss}}$. 反之, 对所有 $n\geq0$, $s^n$ 定义同构 $A^G\xrightarrow{\sim}\Gamma(X,\mathcal O_X(nE))^G$. 每个不变截面 $f$ 的限制等于 $gs^n$, 其中 $g$ 来自 $X\to\Spec A^G$ 下的函数拉回. 因 $X$ 约化, $f=g\cdot s^n$, 从而 $X^{\mathrm{ss}}\subseteq\Spec A$.

现验证 ($\star$). 给定 $\xi$ 与有限阶 $g\in\operatorname{Aut}_{\mathcal X}(\xi_K)$, 对 $[X/G]$ 应用上述结果, 扩张后得到 $\xi':\Spec R\to[X/G]$, 泛点与 $\xi$ 同构, $g$ 延拓为其自同构, 但特殊点未必同构于 $x$. 稳定子群概形 $\operatorname{Stab}_G(X)\subseteq X\times G$ 是闭子概形, 关于 $G$ 在 $X\times G$ 上的乘积作用等变, 其中在自身上作共轭作用. 二元组 $(\xi',g)$ 定义 $\eta:\Spec R\to[\operatorname{Stab}_G(X)/G]$. 将证明扩张后存在 $\eta':\Spec R\to[\operatorname{Stab}_G(\Spec A)/G]$, 且 $\eta'_K\simeq\eta_K$. 类似于~\bookeqref{7.10.26}, 有 $G$-等变图
\[\begin{tikzcd}[ampersand replacement=\&]\operatorname{Stab}_G(\Spec A)\arrow[r,hook]\arrow[dr]\&\operatorname{Stab}_G(X)\arrow[d]\\\&\Spec A^G\times G\end{tikzcd}\]
$\operatorname{Stab}_G(X)$ 在 $\Spec A^G\times G$ 上射影. 断言它关于 $\mathcal O_X(E)$ 的拉回的半稳定轨迹恰为 $\operatorname{Stab}_G(\Spec A)$. 截面 $s$ 拉回为不变截面, 给出一个包含关系. 反之, 若 $(y,h)\in\operatorname{Stab}_G(X)$ 且 $y\notin X^{\mathrm{ss}}=\Spec A$, 对 $y$ 在仿射锥 $\widehat X\to\Spec A^G$ 上的提升 $\widehat y$ 应用 Kempf 最优失稳定理~\bookref{8.6.6}, 得到单参数子群 $\lambda$, 使 $\lim_{t\to0}\lambda(t)\cdot\widehat y$ 存在且在零截面中. 因 $G_y\subseteq P_\lambda$ (习题~\bookref{8.6.13}), $\lim_{t\to0}\lambda(t)\cdot(\widehat y,h)$ 也存在且在 $\operatorname{Stab}_G(X)$ 的仿射锥的零截面中. 所以 $(y,h)$ 不半稳定.

诱导的 GIT 商态射 $\operatorname{Stab}_G(\Spec A)\sslash G\to(\Spec A^G\times G)\sslash G$ 固有, 良模空间 $[\operatorname{Stab}_G(\Spec A)/G]\to\operatorname{Stab}_G(\Spec A)\sslash G$ 普遍闭. 由赋值判据, 扩张后有提升
\[\begin{tikzcd}[ampersand replacement=\&,column sep=small]
\&{[\operatorname{Stab}_G(\Spec A)/G]}\arrow[r]\&\operatorname{Stab}_G(\Spec A)\sslash G\arrow[d]\&\\
\Spec R\arrow[r,"{\eta}"]\arrow[ur,dashed,"{\chi}"]\&\operatorname{Stab}_G(X)\arrow[r]\&\Spec A^G\times G\arrow[r]\&(\Spec A^G\times G)\sslash G
\end{tikzcd}\]
满足 $\chi_K\simeq\xi_K$. 对复合 $\xi':\Spec R\xrightarrow{\chi}[\operatorname{Stab}_G(\Spec A)/G]\to[\Spec A/G]$, 有 $\xi'_K\simeq\xi_K$, 且 $g$ 延拓为 $\xi'$ 的自同构. 为使 $\xi'(0)\simeq x$, 应用下述引理~\bookref{7.10.27}.
\end{proof}
\begin{lemma}\label{src:7.10.27}
设 $\mathcal X=[\Spec A/G]$, 其中 $\Spec A$ 是代数闭域 $\Bbbk$ 上带约化群 $G$ 作用的有限型仿射概形. 设 $\xi,\xi':\Spec R\to\mathcal X$ 来自剩余域为 $\Bbbk$ 的 DVR, 满足 $\xi_K\simeq\xi'_K$, 且 $\xi(0)$ 闭. 对每个 $g\in\operatorname{Aut}_{\mathcal X}(\xi')$, 扩张 $R$ 后存在 $\xi'':\Spec R\to\mathcal X$, 使 $\xi''_K\simeq\xi'_K$, $g|_K$ 延拓为 $\xi''$ 的自同构, 且 $\xi''(0)\simeq\xi(0)$.
\end{lemma}
\begin{proof}
$\xi(0),\xi'(0)$ 在 $\mathcal X\to\Spec A^G$ 的同一纤维中, 所以 $\xi'(0)$ 的闭包含 $\xi(0)$. Kempf 最优失稳定理~\bookref{8.6.5} 给出典范映射 $f:\Theta\to[\Spec A/G]$, 满足 $f(1)\simeq\xi'(0)$, $f(0)\simeq\xi(0)$. 因 $f$ 典范, $f(1)$ 的每个自同构都延拓为映射 $f$ 的自同构. 特别地, $g$ 在 $f(1)=\xi'(0)$ 上的限制唯一延拓为 $f$ 的自同构 $g_f$.

应用奇异粘合引理~\bookref{7.10.28}, 将 $R$ 替换为 $R[\pi^{1/N}]$, 并用 $x\mapsto x^N$ ($N\gg0$) 定义的 $\Theta\to\Theta$ 预复合 $f$ 后, 得到唯一 $\gamma:\phi_R\to\mathcal X$, 满足 $\gamma|_{s=0}\simeq f$, $\gamma|_{t\ne0}\simeq\xi'$. 唯一性保证 $g$ 与 $g_f$ 唯一延拓为 $\gamma$ 的自同构 $g_\gamma$. 所需映射为
\[\xi'':\Spec(R[\sqrt\pi])\xrightarrow{q}\phi_R\xrightarrow{\gamma}\mathcal X,\]
其中 $q$ 在 $(s,t,\pi)$ 坐标中由 $(\sqrt\pi,\sqrt\pi,\pi)$ 定义. 特殊点映到 $0\in\phi_R$. 按构造 $\xi''(0)\simeq\xi(0)$, $\gamma$ 的自同构限制为延拓 $g|_{K(\sqrt\pi)}$ 的 $\xi''$ 的自同构.
\end{proof}
\begin{lemma}[奇异粘合引理]\label{src:7.10.28}
设 $\mathcal X$ 是代数闭域 $\Bbbk$ 上有限型, 具有仿射对角态射的代数叠, $R$ 是剩余域为 $\Bbbk$ 的 DVR. 设 $f:\Theta\to\mathcal X$, $\xi:\Spec R\to\mathcal X$ 配备同构 $f(1)\simeq\xi(0)$. 当 $n\gg0$, 将 $R$ 替换为 $R[\pi^{1/n}]$, 将 $f$ 替换为 $\Theta\xrightarrow{n}\Theta\xrightarrow{f}\mathcal X$ 后, 存在唯一 $\gamma:\phi_R\to\mathcal X$, 使 $\gamma|_{s=0}\simeq f$, $\gamma|_{t\ne0}\simeq\xi$.
\end{lemma}
\begin{proof}
对 $n>0$, 定义 $\phi_R^{n,1}=[\Spec(R[s,t]/(st^n-\pi))/\mathbb G_m]$, 其中 $s,t$ 的权重为 $n,-1$. 有由 $s=0$ 定义的闭浸入 $\Theta\hookrightarrow\phi_R^{n,1}$ 和由 $t\ne0$ 定义的开浸入 $\Spec R\hookrightarrow\phi_R^{n,1}$. 每个 $\phi_R^{n,1}\to\mathcal X$ 限制为 $f:\Theta\to\mathcal X$ 与 $\xi:\Spec R\to\mathcal X$, 带同构 $\xi(0)\simeq f(1)$. 反之将证明 $n\gg0$ 时这样的 $f,\xi$ 典范延拓为 $\phi_R^{n,1}\to\mathcal X$.

令 $C=R[t,\pi/t,\pi/t^2,\ldots]\subseteq R[t]_t$, 图
\[\begin{tikzcd}[ampersand replacement=\&]\Spec\Bbbk[t]_t\arrow[r]\arrow[d]\&\Spec\Bbbk[t]\arrow[d]\\\Spec R[t]_t\arrow[r]\&\Spec C\end{tikzcd}\]
是概形范畴中的推出 (定理~\bookref{B.4.1}). 它 $\mathbb G_m$-等变, 与 $\mathbb G_m$ 作纤维积后仍是推出. 由推出的性质~\bookref{B.4.8}, 对 $\mathbb G_m$ 取商给出具有仿射对角态射的代数叠范畴中的推出方块
\begin{align}\label{src:7.10.29}
\begin{tikzcd}[ampersand replacement=\&]\Spec\Bbbk\arrow[r]\arrow[d]\&\Theta\arrow[d]\arrow[ddr,bend left,"{f}"]\&\\\Spec R\arrow[r]\arrow[drr,bend right,"{\xi}"']\&{[\Spec C/\mathbb G_m]}\arrow[dr,dashed,"{\Psi}"]\&\\\&\&\mathcal X\end{tikzcd}
\end{align}
诱导虚线箭头 $\Psi$. 写 $C=\bigcup C_n$, 其中 $C_n:=R[t,\pi/t^n]\subseteq R[t]_t$. 因 $C_n\cong R[s,t]/(st^n-\pi)$, $[\Spec C_n/\mathbb G_m]\cong\phi_R^{n,1}$. $\mathcal X\to S$ 局部有限呈示, 所以 $n\gg0$ 时 $\Psi$ 唯一分解为 $[\Spec C/\mathbb G_m]\to\phi_R^{n,1}\to\mathcal X$ (习题~\bookref{4.3.33}). 最后用分次代数映射
\[R[s,t]/(st^n-\pi)\to R[\pi^{1/n}][s^{1/n},t]/(s^{1/n}t-\pi),\]
其中 $s^{1/n}$ 权重为 $1$, 所诱导的典范映射 $\phi_{R[\pi^{1/n}]}\to\phi_R^{n,1}$ 预复合唯一的 $\phi_R^{n,1}\to\mathcal X$, 即完成证明.
\end{proof}
\subsection{S-完备性与约化性}\label{src:7.10.6-subsection}
S-完备性刻画分离性 (命题~\bookref{7.9.22}, \bookref{7.9.23}), 蕴含无穿孔惯性 (定理~\bookref{7.10.23}), 从而蕴含保持稳定子的局部商展示的存在 (命题~\bookref{7.10.20}). 现在证明第三个重要性质: 它刻画约化性. 光滑仿射代数群 $G$ 约化当且仅当 $BG$ 为 S-完备, 当且仅当 $G$ 具有 Cartan 分解. 第~\bookref{8.3} 节进一步讨论 Cartan 分解, 并用它证明 Hilbert--Mumford 判据~\bookref{8.4.4}.
\begin{proposition}\label{src:7.10.30}
设 $G$ 是代数闭域 $\Bbbk$ 上的光滑仿射代数群. 以下条件等价:
\begin{enumerate}[(1)]
\item $G$ 约化;
\item $BG$ 是 S-完备的;
\item $G$ 满足\newterm{Cartan 分解}{Cartan decomposition}: 对每个 $\Bbbk$ 上剩余域为 $\Bbbk$, 分式域为 $K$ 的完备 DVR $R$, 以及每个 $g\in G(K)$, 存在 $h_1,h_2\in G(R)$ 和单参数子群 $\lambda:\mathbb G_m\to G$, 使 $g=h_1\lambda|_Kh_2$.
\end{enumerate}
\end{proposition}
\begin{proof}
命题~\bookref{7.9.20} 已给出 (1)$\Rightarrow$(2). 对 (2)$\Rightarrow$(3), 因 $\phi_R\setminus0=\Spec R\cup_{\Spec K}\Spec R$, 元素 $g\in G(K)$ 通过在 $\Spec K$ 上按 $g$ 粘合两个平凡 $G$-挠子, 确定 $\rho_g:\phi_R\setminus0\to BG$. S-完备性给出提升
\begin{align}\label{src:7.10.31}
\begin{tikzcd}[ampersand replacement=\&]\phi_R\setminus0\arrow[r,"{\rho_g}"]\arrow[d]\&BG\\\phi_R\arrow[ur,dashed,"{h}"']\&\end{tikzcd}
\end{align}
限制 $h|_{B\mathbb G_m}$ 对应于单参数子群 $\lambda$ (相差共轭), 它是分解中的候选. 有两个观察:
\begin{itemize}
\item $\rho_g,\rho_{g'}$ 同构当且仅当存在 $h,h'\in G(R)$, 满足 $hg=g'h'$.
\item 对单参数子群 $\lambda$, 复合 $\phi_R\setminus0\hookrightarrow\phi_R\to B\mathbb G_m\xrightarrow{\lambda}BG$ 与 $\rho_{g'}$ 同构, 其中 $g'=\lambda|_K$.
\end{itemize}
所以只需证明~\bookeqref{7.10.31} 的延拓 $h$ 同构于 $\lambda|_{\phi_R}:\phi_R\to B\mathbb G_m\xrightarrow{\lambda}BG$. 令 $\mathcal P,\mathcal P'$ 为分类二者的主 $G$-丛. 因 $G$ 光滑且仿射, $\underline{\operatorname{Isom}}_{\phi_R}(\mathcal P,\mathcal P')\to\phi_R$ 光滑且仿射. 在 $\mathcal X_0:=B\mathbb G_m\hookrightarrow\phi_R$ 上有截面. 对第 $n$ 个幂零加厚 $\mathcal X_n$, 变形理论~\bookref{7.7.12} 与其上同调仿射性给出相容截面. 凝聚 Tannaka 对偶性~\bookref{7.6.9} 及沿 $B\mathbb G_m$ 的凝聚完备性~\bookref{7.6.13} 蕴含
\[\operatorname{Mor}_{\phi_R}(\phi_R,\underline{\operatorname{Isom}}_{\phi_R}(\mathcal P,\mathcal P'))\to\varprojlim_n\operatorname{Mor}_{\phi_R}(\mathcal X_n,\underline{\operatorname{Isom}}_{\phi_R}(\mathcal P,\mathcal P'))\]
是等价, 所以得到全局截面, 即 $\mathcal P\cong\mathcal P'$. (此蕴含的逻辑可逆.)

对 (3)$\Rightarrow$(1), 若 $G$ 非约化, 则幺幂根中有正规子群 $\mathbb G_a\triangleleft R_u(G)$. 因 $G/R_u(G)$ 与 $R_u(G)/\mathbb G_a$ 均仿射, 复合 $B\mathbb G_a\to BR_uG\to BG$ 仿射. 引理~\bookref{7.9.19} 将蕴含 $B\mathbb G_a$ 为 S-完备, 但这是矛盾: 取 $R=\Bbbk[x]$, $K=\Bbbk((x))$, 元素 $x\in\mathbb G_a(K)$ 不能写成 $h_1\lambda|_Kh_2$. 另见 \cite[定理~A]{AHHL21}.
\end{proof}
特别地, 若 $\mathcal X$ 为 S-完备代数叠, 闭点 $x$ 的稳定子 $G_x$ 光滑且仿射, 则 $G_x$ 约化. 在更强的良模空间存在假设下, 习题~\bookref{7.5.27} 已得到此结论.
\subsection{良模空间存在性定理的证明}\label{src:7.10.7-subsection}
两种完备性的必要性已由命题~\bookref{7.9.23} 证明. 现在按第~\bookref{7.10.1-subsection} 节的策略证明充分性.
\begin{proof}[定理~\bookref{7.10.1} 的证明]
因 $\mathcal X$ 为 S-完备且 $\operatorname{char}(\Bbbk)=0$, 每个闭点 $x$ 的稳定子 $G_x$ 线性约化 (命题~\bookref{7.10.30}). 由局部结构定理~\bookref{7.7.1}, 存在仿射平展态射 $f:([\Spec A/G_x],w)\to(\mathcal X,x)$, 在 $x$ 处诱导稳定子群同构. 由两种完备性, 缩小到 $x$ 的开邻域后, 可假设 $f$ 为 $\Theta$-满射且在所有点处保持稳定子 (命题~\bookref{7.10.13}, \bookref{7.10.20}). 因 $\mathcal X$ 拟紧, 可取有限多个闭点 $x_i$ 及这样的 $f_i:[\Spec A_i/G_{x_i}]\to\mathcal X$, 其像覆盖 $\mathcal X$. 选择嵌入 $G_{x_i}\hookrightarrow\mathrm{GL}_n$, 有等价 $[\Spec A_i/G_{x_i}]\cong[(\Spec A_i\times^{G_{x_i}}\mathrm{GL}_n)/\mathrm{GL}_n]$. 置 $A=\prod_i(A_i\times^{G_{x_i}}\mathrm{GL}_n)$, 得到满射, 仿射, 平展态射
\[f:\mathcal X_1:=[\Spec A/\mathrm{GL}_n]\to\mathcal X,\]
它 $\Theta$-满射且在所有点处保持稳定子. 特征为 $0$, 所以有良模空间 $\mathcal X_1\to X_1:=\Spec A^{\mathrm{GL}_N}$. 令 $\mathcal X_2=\mathcal X_1\times_{\mathcal X}\mathcal X_1$. 投影 $p_1,p_2:\mathcal X_2\to X_1$ 也仿射, 平展, $\Theta$-满射且保持稳定子. 因 $f$ 仿射, $\mathcal X_2\cong[\Spec B/\mathrm{GL}_n]$, 有良模空间 $\mathcal X_2\to X_2:=\Spec B^{\mathrm{GL}_n}$. 得到交换图
\begin{align}\label{src:7.10.32}
\begin{tikzcd}[ampersand replacement=\&]\mathcal X_2\arrow[r,shift left,"{p_1}"]\arrow[r,shift right,"{p_2}"']\arrow[d]\&\mathcal X_1\arrow[r,"{f}"]\arrow[d]\&\mathcal X\arrow[d,dashed]\\X_2\arrow[r,shift left,"{q_1}"]\arrow[r,shift right,"{q_2}"']\&X_1\arrow[r,dashed]\&X\end{tikzcd}
\end{align}
由 Luna 基本引理~\bookref{7.5.33}, 左边每个方块都笛卡尔. 良模空间的泛性质 (定理~\bookref{7.5.11}(4)) 将 $\mathcal X_2\rightrightarrows\mathcal X_1$ 的平展群胚结构诱导为 $X_2\rightrightarrows X_1$ 上的平展群胚结构.

断言它是平展等价关系, 即商叠 $[X_1/X_2]$ 是代数空间. 由代数空间的刻画~\bookref{4.6.6}, 只需证明对每个 $\Bbbk$-点 $x_1\in X_1$, $(x_1,x_1)$ 在 $(q_1,q_2):X_2\to X_1\times X_1$ 下有唯一原像. 设 $x_2,x'_2$ 均映到此点, 令 $\widetilde x_2,\widetilde x'_2\in|\mathcal X_2|$ 为其原像中的唯一闭点. 因 $f$ 为 $\Theta$-满射, 它们的四个投影都是 $x_1$ 上的闭点, 故均等于 $x_1$ 上的唯一闭点 $\widetilde x_1$. 因 $f$ 保持稳定子, $\widetilde x_2,\widetilde x'_2$ 的稳定子群与 $\widetilde x_1$ 及其在 $\mathcal X$ 中的像的稳定子相同, 记为 $G$. 所以 $(p_1,p_2):\mathcal X_2\to\mathcal X_1\times\mathcal X_1$ 沿剩余 gerbe $\mathcal G_{(\widetilde x_1,\widetilde x_1)}=BG\times BG\to\mathcal X_1\times\mathcal X_1$ 的纤维积同构于 $BG$, 因而是唯一闭点的剩余 gerbe. 所以 $x_2=x'_2$.

于是 $X=X_1/X_2$ 是代数空间. 由平展下降有态射 $\mathcal X\to X$, 沿 $X_1\to X$ 的拉回为良模空间 $\mathcal X_1\to X_1$. 由良模空间的下降 (引理~\bookref{7.5.24}(2)), $\mathcal X\to X$ 是良模空间. 最后, S-完备性蕴含 $X$ 分离 (命题~\bookref{7.9.23}).
\end{proof}
